\documentclass[12pt]{article}
\usepackage[utf8]{inputenc}
\usepackage[T1]{fontenc}
\usepackage{hyperref}
\title{Tho Language Resources for K-12 Curriculum Development}
\author{Compiled by Claude Sonnet 5 (claude-sonnet-5) for Andr\'as Kornai}
\date{}
\begin{document}
\maketitle

\section{Introduction}
Tho (ISO 639-3: tou; Glottolog: thoo1240), also called Cuoi, is a Vietic (Austroasiatic) language of Nghe An province in Vietnam. Some speakers live in Laos. It is not the Tai-language Tho/Tay. Wikipedia reports about 70,000 speakers, and the Th\h{o} population was 91,430 in 2019. It is closely related to Vietnamese. No school materials were found.

\section{Primary Educational Resources}
None found.

\section{Linguistic and Reference Materials}
\begin{itemize}
\item Ferlus, a paper on Cuoi Cham: \url{https://halshs.archives-ouvertes.fr/halshs-01079521}
\item Wikipedia, Cuoi language: \url{https://en.wikipedia.org/wiki/Cu%E1%BB%91i_language}
\item Wikipedia, Vietic languages: \url{https://en.wikipedia.org/wiki/Vietic_languages}
\item ThoT database, tonal status: \url{https://thot.huma-num.fr/db/thoo1240/tonal_status}
\item Max Planck numerals page for Tho: \url{https://lingweb.eva.mpg.de/channumerals/Tho.htm}
\item Huma-Num language listing for code tou. It cites Nguy\~en H\~uu Ho\`anh (2009), but I could not confirm the title or type: \url{https://cocoon.huma-num.fr/exist/crdo/lexvo.xql?code=tou}
\end{itemize}

\section{Religious and Community Resources}
The Huma-Num listing mentions a New Testament in some documentation, but this was not verified.

\section{Digital and Audio Resources}
None confirmed.

\section{Curriculum Development Context}
No Tho-language school programme was found. Instruction would in practice rely on Vietnamese. Any native-language curriculum would need new documentation.
\end{document}